% Source: https://artofproblemsolving.com/community/c3h2157634_geometry
% Stable identity: aops-topic-2157634-q04
The volume of a cone is one third of the base area times the height.

\begin{asy}
usepackage("amsmath");
size(3cm);
defaultpen(fontsize(10pt));
real r=2, h=12; real k=0.3; defaultpen(fontsize(9pt));
path base=ellipse((0,0),r,r*k);
draw(subpath(base,0,2),dashed); draw(subpath(base,2,4),linewidth(0.9));
draw((-r,0)--(0,h)--(r,0),linewidth(0.9));
draw((0,0)--(0,h),blue+dashed); draw((0,0)--(r,0),red);
draw((-r-0.9,0)--(-r-0.9,h),blue,Arrows(4)); draw((-r-1.2,h)--(-r-0.6,h),blue); draw((-r-1.2,0)--(-r-0.6,0),blue);
label("$12$",(-r-0.9,h/2),W,blue); label("$2$",(r/2,0),S,red);
dot((0,0));
label("\mbox{}",(min(currentpicture,user=true)+max(currentpicture,user=true))/2,black); // keeps the text after the figure black when LaTeX typesets the labels inline
\end{asy}

So the volume is
\[\dfrac{1}{3}\pi\cdot2^2\cdot12=\boxed{16\pi}\approx50.3\text{ cubic metres}.\]
